\begin{table}[H]
\centering\small
\caption{Efecto regional y reescalamiento nacional ilustrativo ante un aumento real de 1\% del mínimo}
\label{tab:chaparro-escala}
\begin{tabularx}{\textwidth}{@{}>{\raggedright\arraybackslash}Xrr@{}}
\toprule
Puestos en un año & Antioquia & Colombia, ilustrativo\\
\midrule
Menor creación de puestos & 44.557 & 246.000 \\
Mayor destrucción de puestos & 5.015 & 28.000 \\
Reducción del saldo neto & 49.572 & 274.000 \\
\bottomrule
\end{tabularx}
\fuente{Fuente: Patiño, Chaparro y coautores, tabla C.1. Cálculo ilustrativo del CJC con afiliados dependientes a cajas de la SSF de marzo de 2025. Los valores nacionales aproximan empleo dependiente afiliado a cajas, se redondean al millar y representan diferencias respecto del escenario sin el aumento. La extrapolación mantiene proporcionalmente los efectos regionales y no constituye una estimación causal para Colombia.}
\end{table}
